\documentclass[11pt]{article}
\usepackage[spanish]{babel}
\usepackage[utf8]{inputenc}
\usepackage{graphicx,array}
\usepackage{amsmath}
\usepackage{multicol}
\usepackage{color}
\usepackage{listings}
\usepackage[inline]{enumitem}
\usepackage[colorlinks]{hyperref}
\usepackage{pdfpages}
\usepackage{amsthm}
\usepackage{physics}
\usepackage{caption}
\usepackage{subcaption}

% Define code style for Python listings
\newtheorem{theorem}{Teorema}
\definecolor{darkgray}{rgb}{0.66, 0.66, 0.66}
\definecolor{darkorange}{rgb}{1.0, 0.55, 0.0}
\definecolor{gray}{rgb}{0.97,0.97,0.99}
\definecolor{teal}{rgb}{0.0, 0.5, 0.5}
\definecolor{comment}{rgb}{0.6, 0, 0.9}

\lstdefinestyle{mystyle}{
    language = Python,
    backgroundcolor=\color{gray},
    commentstyle=\color{comment},
    keywordstyle=\bfseries\color{darkorange},
    numberstyle=\scriptsize\color{darkgray},
    stringstyle=\color{teal},
    basicstyle=\scriptsize\ttfamily,
    breaklines=true,
    captionpos=t,
    numbers=left,
    numbersep=3pt,
    showstringspaces=false,
    tabsize=4
}
\lstset{style=mystyle}

\synctex=1

\title{Modos Atrapados para un Obstáculo Parametrizado en una Guía de Ondas Bidimensional}
\author{Kevin Santiago Sepúlveda García}
\date{\today}

\begin{document}

\maketitle

\begin{abstract}
Este informe presenta la validación numérica del Teorema~2.1 de~\cite{Zhevandrov2025}, que predice la existencia de modos atrapados en una guía de ondas bidimensional con un obstáculo pequeño. Utilizando el Método de Elementos de Frontera (BEM), calculamos los valores propios de la ecuación de Helmholtz para diferentes posiciones del obstáculo y parámetros de perturbación $\varepsilon$. El objetivo principal es determinar el valor máximo de $\varepsilon$, denotado $\varepsilon_{\max}$, para el cual el teorema se cumple. Además, analizamos cómo $\varepsilon_{\max}$ depende de la altura $a$ del obstáculo.
\end{abstract}

\section{Introducción}

Los modos atrapados son oscilaciones localizadas que pueden ocurrir en guías de ondas abiertas (infinitas) debido a la presencia de obstáculos o irregularidades geométricas. Corresponden a soluciones no radiantes de la ecuación de Helmholtz que permanecen confinadas cerca de la perturbación.

La existencia de modos atrapados es de gran interés teórico y práctico, ya que juegan un papel clave en la propagación de ondas, la dispersión y la localización de energía en sistemas físicos como ductos acústicos, canales de agua y guías de ondas cuánticas.

En este informe, validamos el Teorema~2.1 de~\cite{Zhevandrov2025} utilizando simulaciones numéricas basadas en el Método de Elementos de Frontera (BEM). El teorema proporciona una expresión asintótica para el desplazamiento del valor propio asociado con el modo atrapado en función del parámetro $\varepsilon$ y la altura $a$ del obstáculo.

Nuestros objetivos son:
\begin{itemize}
    \item Verificar la expresión analítica para $\sigma$ dada en el teorema.
    \item Determinar el rango de $\varepsilon$ para el cual el teorema permanece válido.
    \item Analizar la dependencia de $\varepsilon_{\max}$ con respecto al parámetro $a$.
\end{itemize}

\section{Marco Teórico}

\subsection{Fenómenos de Propagación de Ondas en Guías Confinadas}

Las guías de ondas son estructuras geométricas que confinan la propagación de ondas en direcciones transversales permitiendo propagación a lo largo de una dirección principal. El fenómeno de confinamiento modal es fundamental en múltiples disciplinas de la ingeniería, desde sistemas de telecomunicaciones hasta acústica aplicada en sistemas de ventilación, motores de aviación y tuberías industriales \cite{Rienstra2015}.

La propagación de ondas en guías confinadas se rige por el principio de modos propios. Para una guía de ondas infinita, el campo de onda puede descomponerse en una suma discreta de modos de propagación, cada uno caracterizado por una relación de dispersión específica que vincula la frecuencia con el número de onda axial. Este conjunto modal forma una base ortonormal que permite expandir cualquier solución como combinación lineal de modos \cite{Dostart2017}.

En sistemas prácticos, la presencia de obstáculos, discontinuidades geométricas o cambios en las propiedades del material introduce complejidad adicional. Estos elementos pueden provocar dispersión modal, acoplamiento entre modos e interacciones complejas que requieren metodologías numéricas avanzadas para ser caracterizadas adecuadamente. Las aplicaciones ingenieriles incluyen el diseño de silenciadores en sistemas de escape de motores de combustión interna, atenuadores acústicos en conductos HVAC, y dispositivos de control de ruido en sistemas de turbomaquinaria \cite{Kirby2005, Rämmal2009}.

\subsection{Guías de Ondas con Obstáculos Rígidos: Fundamentos Matemáticos}

Las \textbf{guías de ondas con obstáculos rígidos} constituyen un problema fundamental en la propagación de ondas con aplicaciones extensas en acústica industrial, electromagnetismo y dinámica de fluidos. Un guía de ondas típica consiste en un canal infinito (o semi-infinito) limitado por paredes rígidas, en cuyo interior se introduce un obstáculo que altera la propagación de las ondas \cite{Linton2008, Pagneux2007, Khallaf2000}.

Matemáticamente, el problema se formula mediante la ecuación de Helmholtz, que es la forma frecuencial de la ecuación de ondas para oscilaciones armónicas en el tiempo. La ecuación de Helmholtz describe el potencial de desplazamiento o presión acústica en un medio homogéneo:

\begin{equation}
-\Delta u = k^2 u, \quad (x, y) \in \Omega,
\end{equation}

donde $k = \omega/c$ es el número de onda, $\omega$ es la frecuencia angular, y $c$ es la velocidad de propagación en el medio.

Los modos de propagación en una guía de ondas se clasifican en dos categorías principales según su relación con la frecuencia de corte: \textbf{modos propagativos} y \textbf{modos evanescentes}. Los modos propagativos tienen números de onda reales y transportan energía a lo largo de la guía. Los modos evanescentes poseen números de onda complejos y decaen exponencialmente a lo largo de la dirección de propagación, sin contribuir al transporte neto de energía a distancias grandes \cite{Pagneux2007}.

Cada modo propagativo posee una \textbf{frecuencia de corte} característica. Para una guía de ancho $2b$ con paredes Dirichlet (rígidas), la $n$-ésima frecuencia de corte es:

\begin{equation}
\Lambda_n = \left( \frac{n \pi}{2b} \right)^2, \quad n = 1, 2, 3, \ldots
\end{equation}

Por debajo de la frecuencia de corte de un modo particular, ese modo es evanescente. Por encima, el modo se propaga. Esta estructura modal es crítica para entender cómo los obstáculos interactúan con el campo de ondas.

\subsection{Modos Atrapados: Concepto y Caracterización}\label{subsec:problem}

Los \textbf{modos atrapados} son soluciones localizadas especiales de la ecuación de Helmholtz que decaen exponencialmente en ambas direcciones a lo largo de la guía de ondas. A diferencia de los modos propagativos ordinarios, los modos atrapados existen solo en frecuencias específicas determinadas por la geometría del obstáculo y las condiciones de frontera \cite{Zhevandrov2025, Khallaf2000}.

Físicamente, un modo atrapado corresponde a una frecuencia a la cual el obstáculo actúa como un resonador que confina la energía ondulatoria localmente. Esta resonancia es posible porque el obstáculo dispersa parte de la energía incidente de forma que se produce interferencia constructiva en la vecindad del obstáculo, generando una oscilación estacionaria localizada.

Matemáticamente, un modo atrapado es un valor propio del operador de Helmholtz (considerado como un problema de valores propios generalizado) que satisface:

\begin{equation}
-\Delta u = k^2 u, \quad (x, y) \in \Omega,
\end{equation}

sujeto a:

\begin{align}
u &= 0, \quad \text{en } \Gamma_\pm = \{y = \pm b\} \quad (\text{paredes de Dirichlet}), \\
\pdv{u}{n} &= 0, \quad \text{en } \gamma \quad (\text{obstáculo de Neumann}), \\
u(x,y) &\to 0 \quad \text{cuando } |x| \to \infty,
\end{align}

con la condición adicional de que $k^2 < \Lambda_1$ (el valor propio está por debajo del primer umbral de propagación).

La existencia de modos atrapados bajo pequeñas perturbaciones es un resultado profundo en la teoría de operadores diferenciales. Evans \emph{et al.} \cite{Evans1994} demostraron que para obstáculos suficientemente simétricos en una guía de ondas, siempre existe al menos un modo atrapado. Este resultado ha sido extendido y refinado por múltiples autores, con contribuciones particularmente importantes de Zhevandrov y colaboradores, que desarrollaron teoría de perturbaciones asintótica para describir cómo emergen modos atrapados cuando un obstáculo pequeño es introducido en una guía de ondas \cite{Zhevandrov2019, Zhevandrov2025}.

\subsection{Métodos Numéricos para Problemas de Valores Propios en Guías de Ondas}

La solución numérica de problemas de valores propios asociados con ecuaciones diferenciales parciales en dominios complejos requiere seleccionar técnicas numéricas apropiadas según las características específicas del problema \cite{Bao2021}.

\subsubsection{Método de Elementos de Frontera (BEM)}

El Método de Elementos de Frontera es particularmente efectivo para dominios no acotados como las guías de ondas infinitas. A diferencia de los métodos de dominio completo (FEM, diferencias finitas), el BEM reduce la dimensionalidad del problema al discretizar únicamente la frontera, resultando en matrices de menor tamaño \cite{Pozrikidis2002}.

Las ventajas principales del BEM para el presente problema son:

\begin{itemize}
\item \textbf{Manejo automático del espacio infinito}: La función de Green fundamental satisface automáticamente la condición de radiación de Sommerfeld en el infinito.

\item \textbf{Reducción dimensional}: Solo se discretiza el contorno del obstáculo ($O(1)$-dimensional para problemas 2D), no todo el dominio físico.

\item \textbf{Alta precisión con malla moderada}: Problemas BEM típicamente requieren fewer degrees of freedom que FEM para la misma precisión.

\item \textbf{Formulación de valores propios clara}: Los valores propios se obtienen como ceros del determinante de la matriz BEM, proporcionando un procedimiento bien definido.

\item \textbf{Uso de funciones de Green simétricas}: Para guías de ondas con paredes Dirichlet, la función de Green simétrica $G_s$ satisface automáticamente las condiciones de Dirichlet en las paredes, simplificando la formulación.
\end{itemize}

La función de Green para una guía de ondas con paredes Dirichlet en $y = \pm b$ se construye mediante el método de imágenes, garantizando que se cumplan las condiciones de frontera en las paredes sin necesidad de discretizar esas fronteras \cite{Linton1992}.

\subsubsection{Método de Coincidencia de Modos (Mode-Matching)}

El método de coincidencia de modos es un enfoque alternativo basado en expandir el campo de onda en términos de modos propios de la guía de ondas no perturbada. En regiones donde la geometría es uniforme, el campo se expresa como:

\begin{equation}
u(x,y) = \sum_{n} A_n e^{i\beta_n x} \psi_n(y),
\end{equation}

donde $\psi_n(y)$ son los modos transversales de la guía sin obstáculo y $\beta_n$ son los números de onda axiales correspondientes.

En discontinuidades o en presencia de obstáculos, se imponen condiciones de coincidencia que requieren continuidad de presión y flujo normal a través de planos de corte. El sistema resultante de ecuaciones lineales proporciona los coeficientes modales y permite determinar los números de onda de resonancia \cite{Klaren1994, Alkuhayli2025}.

El mode-matching es computacionalmente más eficiente que FEM para geometrías 2D simples pero es más restrictivo en cuanto a complejidad geométrica.

\subsubsection{Método de Elementos Finitos (FEM)}

El FEM es una técnica general de discretización aplicable a dominios complejos. Para guías de ondas, se han desarrollado variantes especializadas como el Waveguide FEM (WFEM) que aprovecha la invariancia longitudinal de la estructura \cite{Finnveden2012}.

El WFEM resuelve un problema de valores propios 2D en la sección transversal para obtener directamente los números de onda de propagación en función de la frecuencia. Esta formulación es particularmente eficiente porque reduce el problema 3D a uno 2D, disminuyendo significativamente el costo computacional \cite{Dostart2017}.

Las ventajas del FEM/WFEM incluyen:
\begin{itemize}
\item Capacidad para manejar geometrías altamente complejas.
\item Tratamiento sistemático de condiciones de frontera arbitrarias.
\item Disponibilidad de software maduro y bien validado.
\item Facilidad para incorporar materiales heterogéneos.
\end{itemize}

Sin embargo, para dominios no acotados, es necesario emplear técnicas especiales como capas perfectamente absorbentes (PML) o Dirichlet to Neumman (DtN) para truncar el dominio de cálculo \cite{Bao2021, Casenavea2014}.

\subsection{Aplicaciones Ingenieriles de Modos Atrapados y Resonancias en Guías}

Los fenómenos de resonancia en guías de ondas tienen aplicaciones práticas significativas en ingeniería mecánica y acústica:

\begin{enumerate}
\item \textbf{Silenciadores y atenuadores acústicos}: El diseño de silenciadores para sistemas de escape de motores, conductos HVAC y sistemas de tuberías se basa en el control deliberado de resonancias. Obstáculos estratégicamente posicionados pueden amplificar la dispersión de modos propagativos en bandas de frecuencia específicas, reduciendo la transmisión de sonido \cite{Kirby2005}.

\item \textbf{Filtros acústicos de banda angosta}: Cavidades laterales en guías de ondas actúan como resonadores Helmholtz que atenúan energía en frecuencias de resonancia específicas. Estos dispositivos se utilizan en sistemas de ventilación industrial y aviación \cite{Alkuhayli2025}.

\item \textbf{Sistemas de amortiguamiento en turbomaquinaria}: En turbocompresores y turbinas, los obstáculos en los conductos de admisión y escape pueden crear respuestas resonantes que amplifican o atenúan ciertos componentes del espectro de ruido \cite{Rämmal2009}.

\item \textbf{Cavidades resonantes en sistemas de transmisión de potencia}: En sistemas de tuberías de alta presión, resonancias no deseadas pueden provocar oscilaciones de presión y daño estructural. El diseño cuidadoso de obstáculos internos permite controlar estas oscilaciones \cite{Hein2012}.

\item \textbf{Acoplamiento de modos y Fano resonancias}: En sistemas con múltiples cavidades o discontinuidades, fenómenos de interferencia cuántica (como resonancias Fano) pueden producir picos de reflexión extremadamente afilados, útiles para filtrado de precisión \cite{Kronowetter2023}.

\end{enumerate}

\subsection{Descripción Geométrica del Problema}

Consideramos una guía de ondas bidimensional de ancho $2b$, con paredes ubicadas en $y = \pm b$. El obstáculo es un círculo de radio $r = R_0\varepsilon$ donde $R_0 = 1$ y $\varepsilon \ll 1$ es el parámetro de perturbación. 

Los parámetros geométricos clave son:
\begin{itemize}
    \item $b$: semi-ancho de la guía de ondas (distancia desde el eje central hasta cada pared)
    \item $a$: posición vertical del centro del obstáculo, medida desde la línea central $y=0$ de la guía
    \item $r = \varepsilon$: radio del obstáculo circular
    \item $\alpha = \frac{\pi a}{b}$: posición normalizada del obstáculo
    \item $\varepsilon$: parámetro de perturbación, que mide la magnitud del obstáculo respecto al ancho de la guía. Se asume que $\varepsilon \ll 1$, permitiendo realizar análisis asintóticos y estudiar cómo la presencia de un obstáculo pequeño afecta la existencia de modos atrapados en la guía de ondas.
\end{itemize}

Como se observa en la Figura~\ref{fig:geo}, cuando $a$ aumenta, el centro del obstáculo se acerca a la pared superior ($y=b$). Para valores pequeños de $a$, el obstáculo está cerca del eje central de la guía.

\begin{figure}[h!]
    \centering
    \includegraphics[width=0.6\textwidth]{images/geo_problema.jpeg}
    \caption{Geometría del problema.}
    \label{fig:geo}
\end{figure}

\paragraph{Enunciado del problema.} 
El objetivo de este informe es utilizar el Método de Elementos de Frontera (BEM) para validar el Teorema~2.1 de~\cite{Zhevandrov2025}. Para ello, se busca estudiar el comportamiento asintótico del valor propio asociado al modo atrapado al variar el parámetro de perturbación $\varepsilon$ dentro de un intervalo determinado.

Adicionalmente, se pretende analizar cómo la posición vertical $a$ del obstáculo afecta el valor máximo $\varepsilon_{\max}$ para el cual el modo atrapado sigue existiendo. En otras palabras, se busca cuantificar la relación entre la altura del obstáculo y el rango de valores de $\varepsilon$ que permiten la existencia de modos atrapados.

\subsection{Fundamentos del Método de Elementos de Frontera}

El Método de Elementos de Frontera \cite{Pozrikidis2002} (BEM, por sus siglas en inglés) es una técnica numérica particularmente efectiva para resolver ecuaciones diferenciales parciales en dominios no acotados, como es el caso de las guías de ondas. A diferencia de los métodos de dominio completo como elementos finitos o diferencias finitas, el BEM reduce la dimensionalidad del problema al discretizar únicamente la frontera del dominio, lo que resulta en sistemas de menor tamaño y mayor eficiencia computacional.

Para la ecuación de Helmholtz, el BEM se basa en la representación integral de Green, que expresa la solución en un punto interior del dominio en términos de valores en la frontera. La función de Green $G(P,q)$ satisface
\begin{equation}
\nabla^2 G + k^2 G = -\delta(P-q),
\label{eq:green_basic}
\end{equation}
donde $\delta$ es la distribución delta de Dirac. Para guías de ondas con paredes rígidas (condiciones de Dirichlet), se utiliza la función de Green simétrica $G_s$, que automáticamente satisface las condiciones de frontera en las paredes del canal \cite{Linton1992}.

La ventaja principal del BEM en este contexto es que:
\begin{itemize}
    \item Se consideran automáticamente los efectos del espacio infinito gracias a la función de Green adecuada.
    \item Solo es necesario discretizar el contorno del obstáculo, y no todo el dominio, dado que BEM se basa en la función de Green, la cual representa cómo cada punto de la frontera contribuye al campo en el dominio. En consecuencia, la solución en todo el dominio puede determinarse únicamente a partir de los valores en la frontera.
    \item La precisión del método es alta incluso con un número moderado de elementos de frontera.
    \item El problema de valores propios se reduce a encontrar los valores de $k$ para los cuales la matriz del sistema es singular, porque los valores propios corresponden a las frecuencias en las que existe una solución no trivial. En términos del BEM, esto significa que tras discretizar la frontera se obtiene un sistema lineal de la forma
    \[
    \mathbf{A}(k) \mathbf{u} = \mathbf{0},
    \]
    y $\mathbf{A}(k)$ es singular solo para los valores de $k$ que corresponden a modos atrapados. 
    \item La discretización de la ecuación integral conduce a un sistema algebraico cuyo determinante debe anularse en los valores propios. Esto se demuestra en la subsección \ref{sub:formulation}.

\end{itemize}

\section{Teorema 2.1}

El Teorema~2.1 de~\cite{Zhevandrov2025} proporciona una descripción asintótica del valor propio del modo atrapado para \(\varepsilon\) pequeño.

\begin{theorem}[Teorema 2.1 de~\cite{Zhevandrov2025}]
Para $\varepsilon$ suficientemente pequeño, existe un valor $a = a^* = a_0^* + \mathcal{O}(\varepsilon)$ tal que el problema de Helmholtz \ref{subsec:problem} posee un único modo atrapado en el intervalo $[0, \Lambda_1]$ si $a > a^*$. Si $a \le a^*$, no existen modos atrapados en este intervalo. El valor propio correspondiente es analítico en $\varepsilon$ y $\varepsilon_1 = \varepsilon \ln \varepsilon$, y está dado por $k^2 = \Lambda_1 - \sigma^2$, donde
\begin{align}
\sigma &= \varepsilon^2\frac{\pi^2}{4b^3} \left( \pi \mu\sin^2 \frac{\alpha}{2} - \frac{1}{2} S \cos^2 \frac{\alpha}{2} \right) + \mathcal{O}(\varepsilon^3 \ln \varepsilon), \\
\alpha &= \frac{\pi a}{b}, \\
a_0^* &= \frac{2b}{\pi} \arctan \sqrt{\frac{S}{2\pi\mu}}, \quad \mu = R_0^2, \quad S = \pi R_0^2, \quad r = R_0\varepsilon
\end{align}
\end{theorem}

En este contexto, \(\varepsilon\) es el parámetro de perturbación que mide la magnitud del obstáculo respecto al ancho de la guía de ondas. Se asume que \(\varepsilon \ll 1\), lo que permite aproximar el efecto del obstáculo y analizar cómo su presencia genera un modo atrapado.  

De manera intuitiva, el teorema indica que:
\begin{itemize}
    \item Existe un valor crítico \(a^*\) de la posición vertical del obstáculo por encima del cual se forma un modo atrapado.
    \item El valor propio \(k^2\) del modo atrapado depende de la geometría del obstáculo (\(r\), \(a\)) y de la guía (\(b\)), y se puede expresar asintóticamente en términos de \(\varepsilon\).
    \item La constante \(\mu\) mide la ``fuerza'' del dipolo vertical inducido por el obstáculo, y \(S\) es el área del obstáculo.
\end{itemize}

Para simplificar, en este estudio consideramos \(R_0 = 1\), de modo que \(\mu = 1\), \(S = \pi\) y \(r = \varepsilon\).

\section{Metodología}

\subsection{Formulación del Método de Elementos de Frontera} \label{sub:formulation}

Para calcular numéricamente los números de onda del modo atrapado, empleamos el Método de Elementos de Frontera (BEM) siguiendo el enfoque descrito en~\cite{Linton1992}. 

\paragraph{Formulación del problema.} Recordemos que buscamos soluciones $u(x,y)$ de la ecuación de Helmholtz que satisfagan las condiciones de frontera establecidas en la Sección~\ref{subsec:problem}. Para aplicar el BEM, reformulamos el problema utilizando la representación integral de Green.

\paragraph{Representación integral de la solución en términos de la frontera del obstáculo.} Consideremos un punto $P$ dentro del dominio $\Omega$. Utilizando la fórmula de representación de Green \ref{eq:green_basic}, podemos expresar la solución $u(P)$ en términos de sus valores en la frontera del obstáculo $\gamma$ (ecuacion 3.6 de \cite{Linton1992}):
\begin{equation}
u(P) = \int_{\gamma} \left[ u(q) \pdv{G_s(P,q)}{n_q} - G_s(P,q) \pdv{u}{n}(q) \right] ds_q,
\end{equation}
donde:
\begin{itemize}
    \item $G_s(P,q)$ es la función de Green simétrica para el canal, que satisface automáticamente las condiciones de Dirichlet en las paredes $y = \pm b$
    \item $n_q$ es el vector normal exterior en el punto $q \in \gamma$
    \item $\pdv{u}{n}(q)$ representa la derivada normal de $u$ en la frontera del obstáculo
\end{itemize}

\paragraph{Aplicación de condiciones de frontera.} En el problema \ref{subsec:problem}, el obstáculo satisface condiciones de Neumann homogéneas, es decir, $\pdv{u}{n} = 0$ en $\gamma$. Por lo tanto, la representación integral se simplifica a:
\begin{equation}
u(P) = \int_{\gamma} u(q) \pdv{G_s(P,q)}{n_q} ds_q.
\end{equation}

\paragraph{Ecuación integral de frontera.} Cuando el punto $P$ se aproxima a la frontera $\gamma$ desde el interior del dominio, obtenemos la ecuación integral de frontera:
\begin{equation}
\frac{1}{2} u(p) = \int_{\gamma} u(q) \pdv{G_s(p,q)}{n_q} ds_q, \qquad p \in \gamma.
\end{equation}
El factor $\frac{1}{2}$ surge del límite cuando $P$ tiende a $p$ sobre la frontera.

\paragraph{Discretización.} Para resolver numéricamente esta ecuación integral, parametrizamos la frontera circular del obstáculo en coordenadas polares. Sea $\rho(\theta)$ la representación paramétrica del contorno, con $\theta \in [0, \pi]$ (aprovechando la simetría del problema). Discretizamos el intervalo angular en $M$ puntos equiespaciados:
\begin{equation}
\theta_i = \left(i - \frac{1}{2}\right)\frac{\pi}{M}, \qquad i = 1, 2, \dots, M.
\end{equation}

En cada punto de colocación $\theta_i$, evaluamos la ecuación integral, aproximando la integral mediante una regla de cuadratura:
\begin{equation}
\frac{1}{2} u_i = \frac{\pi}{M} \sum_{j=1}^M u_j K^s_{ij}, \qquad i = 1, \dots, M,
\end{equation}
donde $u_i = u(\theta_i)$ y el núcleo discretizado $K^s_{ij}$ está dado por:
\begin{equation}
K^s_{ij} = 
\begin{cases}
\pdv{G_s(\theta_i, \theta_j)}{n_q}, & \text{si } i \neq j,\\[0.8em]
\pdv{\tilde{G}_s(\theta_i, \theta_i)}{n_q} + \dfrac{\rho_i \rho_i'' - \rho_i^2 - 2\rho_i'^2}{4\pi w_i^3}, & \text{si } i=j.
\end{cases}
\end{equation}
Aquí, $w_i = \sqrt{\rho_i^2 + \rho_i'^2}$. El término diagonal requiere tratamiento especial debido a la singularidad cuando $p = q$, empleando la función de Green regularizada $\tilde{G}_s$.

\paragraph{Sistema algebraico.} Reordenando la ecuación discretizada, obtenemos el sistema matricial homogéneo:
\begin{equation}
\sum_{j=1}^M A_{ij} u_j = 0, \qquad i = 1, \dots, M,
\end{equation}
donde la matriz del sistema está definida por:
\begin{equation}
A_{ij} = \delta_{ij} - \frac{2\pi}{M} K^s_{ij} w_j,
\end{equation}
con $\delta_{ij}$ siendo la delta de Kronecker.

\paragraph{Condición de modo atrapado.} Un modo atrapado corresponde a una solución no trivial del sistema homogéneo, lo cual ocurre cuando la matriz $A$ es singular. Por lo tanto, los valores propios $k$ del problema se obtienen buscando los valores de $k$ que satisfacen:
\begin{equation}
\det(A(k)) = 0.
\end{equation}
Numéricamente, implementamos un algoritmo de búsqueda que varía $k$ en un intervalo de interés y detecta los cruces por cero del determinante.

\subsection{Técnica de Transformación Logarítmica para Análisis de Estabilidad Numérica}

Para mejorar la visualización y estabilidad numérica cuando $\sigma$ varía en varios órdenes de magnitud, introducimos la transformación logarítmica:

\begin{equation}
s = -2 \log_{10}(\sigma) = -\log_{10}(\sigma^2),
\end{equation}

Esta transformación es especialmente útil porque:

\begin{itemize}
\item Comprime el rango dinámico, permitiendo visualizar simultáneamente valores de $\sigma$ que varían en múltiples órdenes de magnitud.

\item Amplifica las diferencias relativas entre valores pequeños de $\sigma$, mejorando la resolución numérica donde las discrepancias entre soluciones analíticas y numéricas son más críticas.

\item Reduce problemas de cancelación catastrófica en aritmética de precisión finita cuando se calculan diferencias de números muy pequeños.

\item Permite una mejor evaluación visual del comportamiento asintótico: desviaciones de la linealidad en un gráfico de $s$ vs $\ln \varepsilon$ indican el inicio de efectos de orden superior no capturados por la aproximación de primer orden.
\end{itemize}

\subsection{Determinación de $\varepsilon_{\max}$}

Para valores fijos de $a$ y $b$, se varía $\varepsilon$ dentro de un rango dado y se resuelve el sistema BEM para cada caso, explorando un intervalo de $k$ en torno a $k_{\text{analytic}}$. Cuando se cumple $\det(A) = 0$, se calcula
\begin{equation}
\sigma_{\text{sol}} = \sqrt{\Lambda_1 - k^2},
\end{equation}
y se compara con la predicción analítica $\sigma_{\text{analytic}}$ del Teorema~2.1.

El teorema se considera válido mientras exista un modo atrapado (es decir, $\sigma_{\text{sol}}$ está definido) y la discrepancia relativa con $\sigma_{\text{analytic}}$ sea pequeña. El mayor valor de $\varepsilon$ que cumple esta condición se denota $\varepsilon_{\max}$ y marca el límite de validez de la aproximación asintótica.

\section{Validación del método BEM}\label{sec:validacion}

Se valida la implementación del método BEM comparando con resultados reportados en \cite{articleMA1}. El objetivo es verificar la precisión del método propio al identificar los valores de $k$ que generan modos atrapados en guías de onda con obstáculos rígidos.

\subsection*{Un obstáculo con condiciones de Dirichlet}

Se considera una guía con un único obstáculo rígido y condición de Dirichlet en el contorno del obstáculo. La Tabla \ref{tab:caso_2_t1} compara los valores de $k$ entre la referencia y la implementación propia. En la Figura \ref{fig:bem_potential} se evidencia la solución $u(x,y)$ para el caso de $r=0.1$ con $k=3.1296$.

\begin{figure}[h!]
    \centering
    \includegraphics[width=1\textwidth]{images/bem_potential.png}
    \caption{Solución numérica de $u(x,y)$ dada por nuestro método BEM para $r=0.1$, donde se obtiene $k=3.1296$.}
    \label{fig:bem_potential}
\end{figure}

\begin{table}[h!]
\centering
\begin{tabular}{|c|c|c|c|}
\hline
$r$ & $k$ (BEM, \cite{articleMA1}) & $k$ (BEM propio) & $k/\pi$ \\ \hline
0.1 & 3.1296 & 3.1296 & 0.996 \\ \hline
0.2 & 3.1415 & 3.1415 & 1.000 \\ \hline
0.3 & 2.9897 & 2.9897 & 0.951 \\ \hline
0.4 & 2.9995 & 2.9995 & 0.954 \\ \hline
0.5 & 3.0717 & 3.0717 & 0.977 \\ \hline
0.6 & 3.1325 & 3.1325 & 0.997 \\ \hline
\end{tabular}
\caption{Comparación de $k$ entre \cite{articleMA1} y la implementación BEM propia para un obstáculo.}
\label{tab:caso_2_t1}
\end{table}

Los resultados muestran excelente concordancia entre la implementación propia y los valores reportados en la literatura, validando así la correcta implementación del método BEM para la detección de modos atrapados en guías de onda con obstáculos rígidos.

\section{Resultados}

\begin{figure}[h!]
    \centering
    \includegraphics[width=0.8\textwidth]{images/plot1.png}
    \caption{Solución numérica $s_{\text{sol}}$ (azul), solución analítica (verde) y error relativo con respecto a la solución analítica $|s_{\text{sol}} - s_{\text{analytic}}|/|s_{\text{analytic}}|$ (naranja) obtenidos usando el BEM.}
    \label{fig:sol}
\end{figure}

\begin{figure}[h!]
    \centering
    \includegraphics[width=0.8\textwidth]{images/plot3.png}
    \caption{Solución numérica para multiples valores de $a$.}
    \label{fig:sol_2}
\end{figure}


\begin{figure}[h!]
    \centering
    \includegraphics[width=0.8\textwidth]{images/plot2.png}
    \caption{Máximo $\varepsilon_{\max}$ en función de la altura $a$ del obstáculo. Al aumentar $a$ se reduce $\varepsilon_{\max}$.}
    \label{fig:sol_a_vs_eps}
\end{figure}

En la Figura~\ref{fig:sol} presentamos la comparación entre los resultados numéricos y analíticos. La curva azul representa la solución numérica $s_{\text{sol}}$ con $a = 0.6$ y $b = 1$, obtenida mediante el Método de Elementos de Frontera (BEM). Se fija $b = 1$ para trabajar en un sistema adimensional, tomando el ancho de la guía como unidad de referencia, lo que simplifica la formulación del problema. La posición vertical $a = 0.6$ se selecciona como un valor intermedio que permite observar claramente la existencia del modo atrapado al variar $\varepsilon$, en la Figura~\ref{fig:sol_2} se muestra el caso para otros valores de $a$. Los puntos donde $s_{\text{sol}} = 0$ indican que no se encontró ningún modo atrapado. La curva naranja muestra el error relativo entre los valores numéricos y analíticos, $s_{\text{sol}}$ y $s_{\text{analytic}}$, respectivamente. El valor máximo para el cual el teorema se cumple es aproximadamente $\varepsilon_{\max} \approx 0.3$. Las soluciones numérica y analítica exhiben fuerte concordancia para valores pequeños de $\varepsilon$. La desviación se vuelve notable para $\varepsilon > 0.3$, sugiriendo que la expansión asintótica derivada en el Teorema~2.1 permanece válida hasta $\varepsilon_{\max} \approx 0.3$.

Adicionalmente, la dependencia del valor crítico $\varepsilon_{\max}$ con respecto a la altura del obstáculo $a$, ilustrada en la Figura~\ref{fig:sol_a_vs_eps}, revela que aumentar $a$ conduce a una disminución en $\varepsilon_{\max}$. Esto implica que posiciones más altas del obstáculo restringen la existencia de modos atrapados, afectando tanto su localización como su estabilidad dentro de la guía de ondas.

\section{Conclusión}

Hemos analizado la existencia de modos atrapados en una guía de ondas con un obstáculo pequeño utilizando tanto métodos numéricos (BEM) como analíticos asintóticos. Los resultados muestran excelente concordancia para valores pequeños de $\varepsilon$, confirmando la validez de la expansión asintótica derivada en el Teorema~2.1 hasta $\varepsilon_{\max} \approx 0.3$.

La dependencia de $\varepsilon_{\max}$ con respecto a la altura $a$ del obstáculo revela que aumentar $a$ reduce el máximo $\varepsilon$ permitido para el cual existen modos atrapados. Esto indica que obstáculos más altos restringen el rango de parámetros que soportan modos atrapados y modifican ligeramente su localización dentro de la guía de ondas.

En general, estos hallazgos proporcionan una guía cuantitativa clara para predecir la existencia y estabilidad de modos atrapados en guías de ondas con diferentes posiciones y tamaños de obstáculos.


\bibliographystyle{unsrt}
\bibliography{references}

\end{document}